% =====================================================================
%  HISTORIAL ACADÈMIC I PROFESSIONAL (HAP) -- Concurs LE-403/729 (UPC)
%  Draft v3 (05/11/2026), English version, prepared from
%  Short_CV_Carlos_Monteiro-2.pdf.
%  Fill in every \TODO{} with real, certifiable data. Only merits
%  certified by 25/11/2026 are assessed (base 3.5).
%  Personal data: ONLY name and surname (base 3.5).
%  Before submitting: uncomment \finaltrue to hide the marks.
%  Output file name: Monteiro_Carlos_HAP.pdf
% =====================================================================
\documentclass[11pt]{article}
\usepackage{lector}
\newcommand{\sectionbreak}{\clearpage} % every section starts on a new page (titlesec)
%\finaltrue
\setupheader{HAP}

\title{\textbf{Historial Acad\`emic i Professional}\\[2pt]\large Academic and Professional Record}
\author{\Large \CandName}
\date{}

\begin{document}
\maketitle
\thispagestyle{plain}
\vspace{0.5cm}
\begin{center}\postdata\end{center}



\newpage
\tableofcontents
\newpage

% =====================================================================
\section{Education and career}

\subsection{Academic qualifications}
\begin{entries}
  \entry{2023--2026}{PhD in Mechanical, Fluids and Aeronautical Engineering}
    {Barcelona East School of Engineering (EEBE), Universitat Polit\`ecnica de Catalunya
    (UPC), \textbf{Spain}}
  \begin{itemize}
    \item Grade: \textbf{Excellent \emph{cum laude}}.\footnote{The thesis was
    completed on 04/09/2026 and the degree issuance fee was paid on 23/09/2026; the degree
    certificate is being issued.} (\certref{phd}{UPC degree fee receipt})
    \item Thesis: \emph{Pseudo-Boiling Wall-Bounded Turbulence in Supercritical
    Flows} (\href{https://hdl.handle.net/2117/473525}{thesis} \docsep see
    Sec.~\ref{sec:thesis})
  \end{itemize}

  \entry{2019--2021}{Master's Degree in Mechanical Engineering}
    {NOVA School of Science and Technology, Universidade NOVA de Lisboa, \textbf{Portugal}}
  \begin{itemize}
    \item Grade: \textbf{16/20}, i.e.\ \textbf{8.17/10} in Spain (\certref{msc-dip}{diploma}
    \docsep \certref{msc}{certificate} \docsep \certref{msc-sup}{Diploma Supplement}
    \docsep \certref{msc-eq}{equivalence})
    \item Thesis (\textbf{19/20}): \emph{Numerical study of a drag reduction system for a
    Formula Student car} (\href{http://hdl.handle.net/10362/149236}{thesis})
  \end{itemize}

  \entry{2019--2020}{Erasmus+ -- Master's courses in Aerospace and Mechanical Engineering}{Politecnico di Milano, \textbf{Italy} (\certref{erasmus}{certificate of attendance})}

  \entry{2016--2019}{Bachelor's Degree in Mechanical Engineering Sciences}
    {NOVA School of Science and Technology, Universidade NOVA de Lisboa, \textbf{Portugal}}
  \begin{itemize}
    \item Grade: \textbf{15/20}, i.e.\ \textbf{7.49/10} in Spain (\certref{bsc}{certificate}
    \docsep \certref{bsc-sup}{Diploma Supplement} \docsep \certref{bsc-eq}{equivalence})
  \end{itemize}
\end{entries}

\subsection{Academic positions}
\begin{entries}
  \entry{01/09/2026 -- present}{Substitute Professor}
    {Department of Fluid Mechanics, Universitat Polit\`ecnica de Catalunya (UPC),
    Barcelona, Spain}
  \begin{itemize}
    \item Teaching: Thermodynamics and Heat Transfer, Fluid Mechanics, Fluid Engineering
    and Fundamentals of Thermal Engineering (\certref{contract-sub}{contract} \docsep see
    Sec.~\ref{sec:teaching})
  \end{itemize}

  \entry{01/09/2026 -- present}{Postdoctoral Researcher}
    {Research group GReCEF, Universitat Polit\`ecnica de Catalunya (UPC), Barcelona, Spain}
  \begin{itemize}
    \item Research: numerical simulation of supercritical turbulent flows, within the ERC
    project SCRAMBLE (\certref{contract-psr}{contract and project addendum} \docsep see
    Sec.~\ref{sec:projects})
  \end{itemize}

  \entry{11/01/2023 -- 31/08/2026}{Predoctoral Researcher}
    {Department of Fluid Mechanics, Universitat Polit\`ecnica de Catalunya (UPC),
    Barcelona, Spain}
  \begin{itemize}
    \item Research: doctoral thesis on supercritical turbulent flows, within the ERC
    project SCRAMBLE (\certref{contract-pre}{contract, extensions and end of contract}
    \docsep see Sec.~\ref{sec:thesis})
  \end{itemize}
\end{entries}

\subsection{Competitive grants and selections}
\begin{entries}
  \entry{2027 \textbf{(scheduled)}}{Mitacs Globalink Research Award}
    {Mitacs, Canada}
  \begin{itemize}
    \item 13-week research stay at the University of British Columbia with Prof.\ M.\
    Rahmani (Section~\ref{sec:stays}).
    \item Status: passed the internal selection of the University of British Columbia.
  \end{itemize}

  \entry{2026}{CTR Summer Program}
    {Center for Turbulence Research (CTR), Stanford University, USA}
  \begin{itemize}
    \item Selected as a participant through the competitive call for research proposals
    (\certref{ctr}{certificate})
  \end{itemize}

  \entry{2023--2026}{\textbf{6$\times$} Supercomputing (HPC) competitive calls}
    {More than 10 million core and GPU hours in total (Section~\ref{sec:hpc})}
  \begin{itemize}
    \item \textbf{5$\times$ Spanish Supercomputing Network (RES)}: MareNostrum~5,
    CTE-AMD (BSC) and \mbox{Finisterrae~III} (CESGA), two of Spain's leading supercomputing centres.
    \item \textbf{1$\times$ Digital Research Alliance of Canada}: Resource Allocation
    Competition (RAC) 2026.
  \end{itemize}

  \entry{2023}{Predoctoral contract}
    {Universitat Polit\`ecnica de Catalunya (UPC), Spain}
  \begin{itemize}
    \item Competitive call 521-729, funded by the ERC project SCRAMBLE (\certref{contract-pre}{contract})
  \end{itemize}
\end{entries}

\subsection[Professional experience]{Professional experience {\normalsize\mdseries(see Sec.~\ref{sec:prof})}}
% The projects and tasks of each position are described in Section~\ref{sec:prof}.
\begin{entries}
  \entry{2021--2022}{Thermal Simulation Engineer}
    {Continental Engineering Services, Porto, Portugal (\certref{continental}{certificate})}

  \entry{2019}{Design Engineering Intern\nocert}
    {Cubot\'onic -- Ind\'ustria Metalomec\^anica, Malveira, Portugal}
\end{entries}

\subsection[Other experience]{Other experience {\normalsize\mdseries(see Sec.~\ref{sec:prof})}}
% The tasks and results are described in Section~\ref{sec:prof}.
\begin{entries}
  \entry{2020--2021}{FST Lisboa Formula Student team -- Aerodynamics and CFD Engineer\nocert}
    {Universidade de Lisboa, Portugal
    (\href{http://hdl.handle.net/10362/149236}{Master's thesis})}
\end{entries}

\subsection{Continuing education}\label{sec:continuing}
Scientific and technical schools and courses, most of them taken on my own initiative.
Teacher training is presented separately in Section~\ref{sec:training}.
\begin{itemize}[itemsep=1pt, parsep=0pt]
  \item \course{2024}{Compute Ontario Summer School}{\textbf{Digital Research Alliance of
  Canada}\nocert}{}
  \item \course{2023}{Summer School on Computational Fluid Dynamics \& SuperComputing}
  {\textbf{Gran Sasso Science Institute}, Italy~(16--22/07/2023)\nocert}{\certref{cfdpar}{letter of admission}}
  \item \course{2023}{Hands on Machine Learning for Fluid Dynamics}{\textbf{von Karman
  Institute}, Belgium~(30~h)}{\certref{vki}{certificate}}
  \item \course{2023}{Machine Learning for Mechanical Engineers}{\textbf{Flowthermolab},
  online~(40~h)}{\certref{ml-ftl}{certificate}}
  \item \course{2022}{Python for Engineers}{\textbf{Flowthermolab}, online~(20~h)}
  {\certref{python}{certificate}}
  \item \course{2022}{CFD Turbulence Modeling: Beyond the Basics}{\textbf{Ansys}~(4~h)}{\certref{ansys}{certificate}}
  \item \course{2022}{CFD Turbulence Modeling: Scale Resolving Simulations}{\textbf{Ansys}~(2.5~h)}{\certref{ansys-srs}{certificate}}
  \item \course{2021}{Simscape Onramp}{\textbf{MathWorks}, self-paced}
  {\certref{simscape}{certificate}}
  \item \course{2021}{Simulink Onramp}{\textbf{MathWorks}, self-paced}
  {\certref{simulink}{certificate}}
  \item \course{2021}{Computational Fluid Dynamics Fundamentals Course 3}{\textbf{Udemy}~(3~h)}
  {\certref{udemy3}{certificate}}
  \item \course{2020}{Workshop de Introdu\c{c}\~ao ao ANSYS}{\textbf{ESSS}~(8~h)}
  {\certref{esss}{certificate}}
  \item \course{2020}{Computational Fluid Dynamics Fundamentals Course 2}{\textbf{Udemy}~(2.5~h)}
  {\certref{udemy}{certificate}}
  \item \course{2020}{Computational Fluid Dynamics Fundamentals Course}{\textbf{Udemy}~(2.5~h)}
  {\certref{udemy1}{certificate}}
  \item \course{2019}{FSG Academy -- Aerodynamics}{\textbf{Formula Student Germany}\nocert}{}
  \item \course{2019}{Vehicle aerodynamics and wind-tunnel testing}{\textbf{Politecnico di Milano}\nocert}{}
\end{itemize}

% =====================================================================
\section{Teaching record and teacher training}\label{sec:teaching}
% \NOTE{Criterion 1 (up to 10 points). Attach the official UPC teaching certificate
% with hours per subject and semester. The 2026--27 hours must be
% certified by the department.}

\subsection{Teaching in official Bachelor's degrees (UPC)}
Every subject I teach (Thermodynamics and Heat Transfer, Fluid Mechanics, Fluid Engineering
and Thermal Engineering) belongs to the profile of the position. My lectures
combine real heat-transfer and fluid-flow cases from industrial practice with current
questions from my doctoral research on real-fluid thermodynamics and turbulence. Hours are
class hours, as certified by the UPC.

\begin{center}\small
\begin{tabularx}{\textwidth}{@{}>{\raggedright\arraybackslash}p{5.0cm}p{1.35cm}>{\raggedright\arraybackslash}Xp{2.0cm}r@{}}
\toprule
\textbf{Subject (UPC code)} & \textbf{School} & \textbf{Activity (groups)} & \textbf{Year, sem.} & \textbf{Hours} \\
\midrule
Thermodynamics and Heat Transfer (820020, TTC) & EEBE & Full course (M1) & 2026--27, S1 & 45 \\
Fluid Mechanics (820013, MF) & EEBE & Laboratory (M33, M34) & 2026--27, S1 & 30 \\
Fundamentals of Thermal Engineering (340038, FENT-F) & EPSEVG & Laboratory (N3111, N3112) & 2026--27, S1 & 15 \\
Fluid Mechanics (820013, MF) & EEBE & Laboratory (M35, T25) & 2025--26, S1 & 30 \\
Fluid Mechanics (820013, MF) & EEBE & Laboratory (M33, M34) & 2025--26, S2 & 30 \\
Fluid Mechanics (820013, MF) & EEBE & Laboratory (T22, T24) & 2024--25, S2 & 30 \\
\midrule
\multicolumn{4}{@{}r}{\textbf{Total, 2024--25 to 2026--27 S1}} & \textbf{180} \\
\midrule
\multicolumn{5}{@{}l}{\emph{Assigned for 2026--27, S2 (spring 2027)} \textbf{(scheduled)}}\\
Thermodynamics and Heat Transfer (820020, TTC) & EEBE & Theory, one further group & 2026--27, S2 & 45 \\
Fluid Engineering (\emph{Enginyeria de Fluids}, 340058, ENFL) & EPSEVG & Theory, groups M65 and M66\textsuperscript{a} & 2026--27, S2 & 92 \\
Fluid Mechanics (820013, MF) & EEBE & Laboratory, one group & 2026--27, S2 & 15 \\
\midrule
\multicolumn{4}{@{}r}{\textbf{Total 2026--27 (90 h in S1 + 152 h in S2)}\textsuperscript{b}} & \textbf{242} \\
\bottomrule
\end{tabularx}
\end{center}
{\footnotesize \textsuperscript{a}46 sessions of 2 h; one weekly session per group is held only in A and B1 weeks (\href{https://drive.google.com/file/d/1TGufsft81vD3qv27ZWpKPmW_KSVtG-hR/view}{see calendar}).
\\\textsuperscript{b}Full contract load (up to 240 teaching hours); the extra 2 h were agreed with the department.\par}
\certlink{teaching}{see UPC certificate, 2024--26} \quad
\certlink{assignment}{see 2026--27 assignment}

\subsection{Course coordination and assessment}
\begin{itemize}
  \item \textbf{2026--27} -- \textbf{Thermodynamics and Heat Transfer} (EEBE). Full course
  (6 ECTS) of group M1, and support to the coordination: exam preparation and
  organisation with Julia-based tools, marking (exams and continuous assessment) and exam
  review
  \certlink{ttc-coord}{see statement}.
  \item \textbf{2026--27 (scheduled)} -- \textbf{Fluid Engineering} (EPSEVG). Course
  coordinator and full course (6 ECTS) of the theory groups M65 and M66, spring 2027.
\end{itemize}

\subsection{Supervision and assessment of academic work}
\begin{itemize}
  \item \textbf{Member of a TFG assessment committee} (EPSEVG), 01--03/02/2027
  \textbf{(scheduled)} -- Bachelor's thesis defences, 2026--27 autumn call
  (\href{https://drive.google.com/file/d/1TGufsft81vD3qv27ZWpKPmW_KSVtG-hR/view}{see academic calendar}).
  \item \textbf{Academic tutor of a curricular internship} of Mag\'i Boada Sivill
  (Bachelor's Degree in Mech.\ Eng.), 2026--27 \textbf{(ongoing)} -- Data analysis of
  extreme events in turbulent flows \certlink{internship}{see project}.
  \item \textbf{Co-supervisor of a Bachelor's thesis (TFG)} of J\'ulia Aguilera Jofre
  (Bachelor's Degree in Mech.\ Eng.), 2023 -- \emph{Microconfined high-pressure
  supercritical turbulent square duct flow}
  (\href{https://hdl.handle.net/2117/395741}{UPCommons})
  \certlink{teaching}{see UPC certificate, p.~3}.
\end{itemize}

% \subsection{Teaching material and innovation}\label{sec:innovation}
% \begin{itemize}
%   \item \textbf{2026--27} -- \textbf{Automated preparation of examinations in Julia} for
%   Thermodynamics and Heat Transfer (EEBE) \certlink{ttc-coord}{see statement}.
%   \item \textbf{2026--27} -- \textbf{Organisation of the TTC theory course in Atenea}, the
%   UPC virtual campus, applying what I learnt in the ICE course \emph{M2.4 -- Atenea: Curs
%   de nivell mitj\`a} (conditional activities, groups and attendance tracking).
%   \item \textbf{2024--2026} -- \textbf{Feedback-based assessment in the Fluid Mechanics
%   laboratories} (EEBE). I return individual written comments on every laboratory report,
%   so that students improve from one report to the next; the students highlight this in
%   the surveys (Section~\ref{sec:surveys}).
% \end{itemize}

\subsection{Teaching quality indicators}\label{sec:surveys}
The official UPC student surveys (\emph{Enquestes de doc\`encia}, scale 1--5) give an
overall mean of \textbf{4.42} for my Fluid Mechanics (820013) laboratory groups in
2025--26, as certified by the UPC. The table shows the groups of 2025--26 S2, which rated
my teaching well above the department, school and university averages.
\certlink{teaching}{see UPC certificate, p.~5}
\certlink{surveys}{see detailed results and student comments}
\begin{center}\small
\begin{tabular}{@{}lccccc@{}}
\toprule
 & \multicolumn{2}{c}{Group} & \multicolumn{3}{c}{Averages, 2025--26 S2} \\
\cmidrule(lr){2-3}\cmidrule(l){4-6}
Question & M33 & M34 & Dept. & EEBE & UPC \\
\midrule
Clear contents; resolves doubts & \textbf{4.80} & \textbf{4.80} & 4.04 & 4.12 & 4.07 \\
Has helped me learn & \textbf{4.80} & \textbf{4.60} & 3.91 & 4.03 & 3.99 \\
\midrule
Responses (participation) & 10/10 (100\%) & 5/6 (83\%) & & & \\
\bottomrule
\end{tabular}
\end{center}
Selected student comments (2025--26, S2, in the original language):
\begin{itemize}
  \item \emph{``Molt bon professor. Explica d'una manera clara i correcte els conceptes i en
  cada informe t'ajuda a millorar progressivament gr\`acies a les seves correcions. Fa que
  tinguis inter\`es per l'assignatura!''}
  \item \emph{``Molt bon professor, ja que t'ajuda en tot i es preocupa de que tothom ho hagi
  ent\`es.''}
\end{itemize}

\subsection{Teaching training}\label{sec:training}
% \NOTE{Criterion 3 (up to 5 points). Only courses completed and certified by 25/11/2026 count.
% The ICE certificate requires a positive assessment of the proposed assignment: submit it and
% request the certificate as soon as the course ends. A single 3 h course scores little; it is
% advisable to add other ICE courses that finish before 25/11.}
\subsubsection*{Courses certified by the UPC}
\begin{itemize}[itemsep=1pt]
  \item \textbf{23--25/11/2026 (enrolled)} -- \emph{M2.4 -- Atenea: Curs de nivell avan\c{c}ats}, ICE--UPC.
  \item \textbf{2026} -- \emph{M2.4 -- Atenea: Curs de nivell mitj\`a}, ICE--UPC \certlink{ice1}{certificate}.
  \item \textbf{09/12/2026 (scheduled)} -- \emph{Taller sobre gesti\'o dels conflictes}, Mediaci\'o UPC (ICE-recognised).
\end{itemize}

\subsubsection*{Other teacher-training courses}
\begin{itemize}[itemsep=1pt]
  \item \textbf{2026} -- \emph{Introduction to Gemini for Education}, Google for Education
  (\href{https://edu.google.accredible.com/9e5d562b-d731-4909-a64e-02fec77496c1}{credential}).
  \item \textbf{2026} -- \emph{Create a lesson plan and supporting materials}, Google for Education
  (\href{https://edu.google.accredible.com/40c5d590-11e2-4418-a7be-5a64a7268e6c}{credential}).
  \item \textbf{2026} -- \emph{Generative AI for Educators and Teachers -- Monsha Certified AI
  Educator: Level 1}, Monsha \certlink{monsha}{certificate}.
\end{itemize}

\subsection{Languages}\label{sec:languages}
\begin{entries}
  \entry{Native}{Portuguese}{}
  \entry{Proficient}{English}{Working language of my research (Sections~\ref{sec:publications} and~\ref{sec:stays}); no formal certificate.}
  \entry{Proficient}{Spanish}{Language of instruction of all my teaching at the UPC (Section~\ref{sec:teaching}); no formal certificate.}
  \entry{A2}{Catalan}{\emph{Catal\`a B\`asic 1} and \emph{3} (A2.1, A2.2; 100 h), SLT--UPC, 2025
  (\certref{catalan1}{A2.1} \docsep \certref{catalan2}{A2.2}); B1 in 2026 \textbf{(scheduled)}.}
\end{entries}

% =====================================================================
\section{Research record}
% \NOTE{Criterion 2 (up to 10 points).}
I have been a member of the research group
\href{https://mf.upc.edu/ca/recerca-1/grecef-grup-de-recerca-en-ciencia-i-enginyeria-de-fluids}{GReCEF
(\emph{Grup de Recerca en Ci\`encia i Enginyeria de Fluids})} of the Department of Fluid
Mechanics, UPC, since January 2023.

\subsection{Researcher identifiers and indicators}
\begin{itemize}
  \item Identifiers: ORCID \href{https://orcid.org/0009-0004-0675-6113}{0009-0004-0675-6113}
  \docsep \href{https://scholar.google.com/citations?user=5nNJuHsAAAAJ&hl=en}{Google Scholar}
  \docsep \href{https://www.researchgate.net/profile/Carlos-Monteiro-19}{ResearchGate} \docsep
  \href{https://futur.upc.edu/CarlosAlexandreBandarrinhaMonteiro}{FUTUR (UPC)}
  \item Publications: 7 articles in indexed journals (2024--2026), 4 articles submitted (2 in the \emph{Journal of Fluid Mechanics}, 1 in \emph{SoftwareX} and
  1 in \emph{Scientific Data}), 2 conference papers and 2 open datasets.
  \item Authorship: first author of 4 of the 7 journal articles, and first or second
  author of all of them.
  \item Quartiles: all 7 journal articles published in Q1 journals.
  \item Citations: 81, h-index 5, i10-index 4 (Google Scholar, 11/11/2026), growing every
  year (1 in 2023, 14 in 2024, 24 in 2025 and 42 in 2026).
  \item Recognition: \emph{Editor's Pick} of \emph{Physics of Fluids} (2026) for [\hyperlink{pub-J2}{J2}].
\end{itemize}
% \NOTE{quartiles to be checked against SJR/JCR before submission}

\subsection{Doctoral thesis}\label{sec:thesis}
\begin{tabular}{@{}p{2.6cm}p{13.3cm}@{}}
Title:       & \emph{Pseudo-Boiling Wall-Bounded Turbulence in Supercritical Flows: Direct
               Numerical Simulation \& Flow Physics} \\
Centre:      & Universitat Polit\`ecnica de Catalunya (UPC), Department of Fluid Mechanics \\
Programme:   & PhD in Mechanical, Fluids and Aeronautical Engineering \\
Supervisor:  & Dr.\ Llu\'is Jofre Cruanyes \\
Completed:   & 04/09/2026 \\
Grade:       & \textbf{Excellent \emph{cum laude}} \\
Repository:  & \href{https://hdl.handle.net/2117/473525}{hdl.handle.net/2117/473525} (UPCommons) \\
\end{tabular}

\medskip
The thesis develops a high-fidelity computational framework for the direct numerical
simulation (DNS) of high-pressure transcritical turbulent flows and uses it to explain
their physics. Its main contributions are:
\begin{itemize}
  \item \textbf{Computational framework.} I redesigned the multiphysics solver RHEA into
  RHEA-X, a GPU-resident version based on MPI and OpenACC, which shows near-ideal strong
  and weak scaling up to large numbers of GPUs [\hyperlink{pub-J1}{J1}].
  \item \textbf{Resolution requirements.} I showed that the classical grid-resolution
  criteria, developed for incompressible flows, are not sufficient near the pseudo-boiling
  region, and proposed new resolution guidelines for DNS of high-pressure transcritical
  wall turbulence [\hyperlink{pub-J5}{J5}].
  \item \textbf{Flow physics.} I showed that microconfined transcritical duct flows can
  move from laminar to turbulence-like states, driven by the density gradients of
  pseudo-boiling, with stronger secondary motions and enhanced heat transfer [\hyperlink{pub-J4}{J4}]; and
  that buoyancy acts mainly by modulating the secondary flows, which redistribute momentum
  and restructure the turbulence in transcritical ducts [\hyperlink{pub-J2}{J2}].
\end{itemize}

\subsection{Research lines}\label{sec:lines}
\begin{revise}
My research is organised in four lines that together cover the three research lines of the
position: (a) multiphysics turbulent flows, compressible and supercritical fluids;
(b) data-driven prediction of engineering flows, extreme events and causality; and
(c) modelling and computation with distributed (CPU) and accelerated (GPU) programming.
The publications supporting each line (Section~\ref{sec:publications}) are cited in brackets.
\begin{enumerate}[label=\textbf{\arabic*.}, leftmargin=*, itemsep=4pt]
  \item \textbf{Supercritical fluid flows} (profile line a). Direct numerical simulation
  (DNS) of pseudo-boiling wall-bounded turbulence with complex real-fluid thermodynamics.
  Core of the PhD thesis.
  \begin{itemize}
    \item DNS methodology: resolution requirements and standards [\hyperlink{pub-J5}{J5}], and minimal
    flow units [\hyperlink{pub-J6}{J6}].
    \item Laminar, transitional and turbulent regimes of microconfined channel and duct
    flows [\hyperlink{pub-J3}{J3}, \hyperlink{pub-J4}{J4}, \hyperlink{pub-J7}{J7}].
  \end{itemize}
  \item \textbf{Multiphysics turbulence in engineering flows} (profile lines a and b).
  High-fidelity analysis of turbulence across applications and flow speeds.
  \begin{itemize}
    \item Complex geometries: square ducts, secondary flows and anisotropic turbulence, and
    their coupling with buoyancy and thermodynamics [\hyperlink{pub-J2}{J2}, \hyperlink{pub-J4}{J4}, \hyperlink{pub-R1}{R1}, \hyperlink{pub-C1}{C1}].
    \item Environmental flows: fundamentals of turbulence and oceanic flows, including the
    aggregation of microplastics (with UBC).
    \item External aerodynamics at different flow speeds: from racing-car aerodynamics
    (Formula Student, drag-reduction system, Master's thesis) [\hyperlink{pub-C2}{C2}] to DNS databases of
    airfoil flows (NACA~0012) [\hyperlink{pub-R5}{R5}].
  \end{itemize}
  \item \textbf{Causality, information theory and data-driven methods} (profile line b).
  Cause-and-effect analysis and prediction of turbulent flows, with emphasis on extreme
  events.
  \begin{itemize}
    \item Causal analysis with transfer entropy, non-stationary transfer entropy and
    Lagrangian information-theoretic tools, applied to extreme wall-heat-flux events in
    supercritical and hypersonic flows; started with the CTR 2026 Summer Program and
    continued with the RES allocation IM-2026-1-0034 and the collaborations of
    Section~\ref{sec:collab} [\hyperlink{pub-R2}{R2}, \hyperlink{pub-R3}{R3}].
    \item Compact neural networks to collapse and discover wall-normal profiles, trained on
    the SCRAMBLE supercritical-flow datasets.
    \item Reduced-order (SVD-based) surrogate models and lookup tables to replace expensive
    real-fluid thermodynamic models [\hyperlink{pub-R4}{R4}].
    \item Flow reconstruction from sparse sensors (RES IM-2025-1/2-0017), data-informed
    analysis of experiments [\hyperlink{pub-J7}{J7}], and open databases for machine learning
    [\hyperlink{pub-J3}{J3}, \hyperlink{pub-D1}{D1}, \hyperlink{pub-D2}{D2}].
  \end{itemize}
  \item \textbf{GPU acceleration of high-fidelity solvers} (profile line c).
  \begin{itemize}
    \item Porting of legacy multiphysics compressible-turbulence DNS codes to GPUs with
    pragma-based paradigms (OpenACC) and MPI, and kernel-level optimisation: GPU
    acceleration of the open-source solver RHEA [\hyperlink{pub-J1}{J1}].
    \item BSC MareNostrum Hackathons (2024 and 2026) and Open Accelerated Computing Summit
    (2023); large-scale simulations on MareNostrum~5, Finisterrae~III and Canadian
    national systems (Section~\ref{sec:hpc}).
  \end{itemize}
\end{enumerate}
\end{revise}

\subsection{Publications}\label{sec:publications}
\subsubsection*{Articles in indexed journals}
\begin{enumerate}[label={[J\arabic*]}]
  \item \hypertarget{pub-J1}{}\textbf{C.\ Monteiro}, X.\ \'Alvarez-Farr\'e, D.\ Mart\'in, F.\ Pizzi, L.\ Jofre.
  RHEA-X: Distributed GPU-resident high-fidelity flow solver for multiphysics turbulence.
  \emph{Computer Physics Communications}, 110233, 2026.
  \href{https://doi.org/10.1016/j.cpc.2026.110233}{doi:10.1016/j.cpc.2026.110233} [Q1].
  \item \hypertarget{pub-J2}{}\textbf{C.\ Monteiro}, G.\ Barea, L.\ Jofre. Buoyancy-driven secondary flows and
  momentum redistribution in high-pressure transcritical duct turbulence.
  \emph{Physics of Fluids}, 065127, 2026.
  \href{https://doi.org/10.1063/5.0337547}{doi:10.1063/5.0337547} [Q1]. \textbf{Editor's Pick.}
  \item \hypertarget{pub-J3}{}A.\ Abdellatif, \textbf{C.\ Monteiro}, M.\ Bernades, L.\ Jofre. Microconfined
  high-pressure transcritical channel flow database: laminar, transitional \& turbulent
  regimes. \emph{Scientific Data}, 11(1), 2024.
  \href{https://doi.org/10.1038/s41597-024-03757-4}{doi:10.1038/s41597-024-03757-4} [Q1].
  \item \hypertarget{pub-J4}{}\textbf{C.\ Monteiro}, L.\ Jofre. Flow regime analysis of high-pressure
  transcritical fluids in microducts. \emph{International Journal of Heat and Mass
  Transfer}, 224, 125295, 2024.
  \href{https://doi.org/10.1016/j.ijheatmasstransfer.2024.125295}{doi:10.1016/j.ijheatmasstransfer.2024.125295} [Q1].
  \item \hypertarget{pub-J5}{}\textbf{C.\ Monteiro}, L.\ Jofre. Resolution standards for direct numerical
  simulation of wall turbulence in high-pressure transcritical fluids. \emph{Physics of
  Fluids}, 36(12), 2024.
  \href{https://doi.org/10.1063/5.0244472}{doi:10.1063/5.0244472} [Q1].
  \item \hypertarget{pub-J6}{}R.\ El Mansy, \textbf{C.\ Monteiro}, F.\ Mellibovsky, L.\ Jofre. Minimal flow unit
  of wall-bounded high-pressure transcritical turbulence. \emph{Physics of Fluids},
  36(12), 2024.
  \href{https://doi.org/10.1063/5.0243832}{doi:10.1063/5.0243832} [Q1].
  \item \hypertarget{pub-J7}{}E.\ Hurt\'an, \textbf{C.\ Monteiro}, M.\ Jofre, J.\ Casals-Terr\'e, L.\ Jofre.
  Data-informed characterization of spatio-temporal scales in experiments of microconfined
  high-pressure transcritical turbulence. \emph{Experimental Thermal and Fluid Science},
  159, 111282, 2024.
  \href{https://doi.org/10.1016/j.expthermflusci.2024.111282}{doi:10.1016/j.expthermflusci.2024.111282} [Q1].
\end{enumerate}

\subsubsection*{Submitted articles and proceedings}
\begin{enumerate}[label={[R\arabic*]}]
  \item \hypertarget{pub-R1}{}\textbf{C.\ Monteiro}, L.\ Jofre. Turbulent supercritical fluid flow in
  square ducts. \emph{Journal of Fluid Mechanics}, submitted.
  \item \hypertarget{pub-R2}{}\textbf{C.\ Monteiro}, D.\ Mart\'in, F.\ Pizzi, A.\ Chandra, L.\ Jofre. Causal structure of extreme
  wall-heat-flux events in supercritical turbulence. \emph{Journal of Fluid Mechanics},
  submitted.
  \item \hypertarget{pub-R3}{}\textbf{C.\ Monteiro}, D.\ Mart\'in, F.\ Pizzi, A.\ Chandra, L.\ Jofre.
  Causal analysis of extreme wall heat flux mechanisms in supercritical turbulence.
  \emph{Center for Turbulence Research, Proceedings of the Summer Program 2026}, Stanford
  University, to appear.
  \item \hypertarget{pub-R4}{}T.\ Pinto, O.\ Cl\'ery-R\'emy, \textbf{C.\ Monteiro}, L.\ Jofre.
  Thermophysical surrogate modeling software. \emph{SoftwareX}, submitted. Open-source
  Python tool that compresses tabulated real-fluid properties into low-rank (SVD)
  surrogates, used in the solver RHEA for DNS of supercritical CO$_2$ turbulence.
  \item \hypertarget{pub-R5}{}E.\ Mart\'in, \textbf{C.\ Monteiro}, D.\ Mart\'in, L.\ Jofre,
  J.\ Calafell. NACA0012 database. \emph{Scientific Data}, submitted. Direct numerical
  simulation database of the transitional and separated flow around a NACA~0012 airfoil at
  low to moderate Reynolds numbers.
\end{enumerate}

\subsubsection*{Conference papers}
\begin{enumerate}[label={[C\arabic*]}]
  \item \hypertarget{pub-C1}{}\textbf{C.\ Monteiro}, L.\ Jofre. Buoyant high-pressure transcritical fluid duct
  flows. \emph{9th European Congress on Computational Methods in Applied Sciences and
  Engineering (ECCOMAS 2024)}, Lisbon, 2024.
  \href{https://doi.org/10.23967/eccomas.2024.054}{doi:10.23967/eccomas.2024.054}. Oral presentation.
  \certlink{eccomas}{see attendance certificate}
  \item \hypertarget{pub-C2}{}\textbf{C.\ Monteiro}, M.\ Brito, D.\ Vieira. Numerical iterative analysis of drag
  reduction system for a racing car rear wing. \emph{SAE Technical Paper Series}, 2023.
  \href{https://doi.org/10.4271/2023-01-5079}{doi:10.4271/2023-01-5079}.
\end{enumerate}

\subsubsection*{Datasets}
\begin{enumerate}[label={[D\arabic*]}]
  \item \hypertarget{pub-D1}{}W.\ T.\ Chung, M.\ Ihme, K.\ S.\ Jung, \emph{et al.} (incl.\ \textbf{C.\ Monteiro}).
  BLASTNet multiphysics simulation dataset, v3.0. Zenodo, 2026.
  \href{https://doi.org/10.5281/zenodo.21348075}{doi:10.5281/zenodo.21348075}.
  \item \hypertarget{pub-D2}{}A.\ Abdellatif, \textbf{C.\ Monteiro}, L.\ Jofre. Microconfined high-pressure
  transcritical channel flow dataset, v3. Figshare+, 2023.
  \href{https://doi.org/10.25452/figshare.plus.24486904.v3}{doi:10.25452/figshare.plus.24486904.v3}.
\end{enumerate}

\subsection{Invited talks and presentations}
\begin{itemize}
  \item \textbf{30/08--03/09/2027 (scheduled)} -- Presentation at the
  \href{https://efdc3.com/}{3rd European Fluid Dynamics Conference (EFDC3)}, Zurich
  (Switzerland): \emph{Turbulent square duct flows of supercritical fluids}.
  \item \textbf{16--19/06/2027 (scheduled)} -- Presentation at the
  \href{https://thmt-27.org/}{12th International Symposium on Turbulence, Heat and Mass
  Transfer (THMT-27)}, Niagara Falls (Canada): \emph{Causal structure of extreme heat flux
  events in supercritical flows}.
  \item \textbf{2026} -- Presentation of the project \emph{Hunting Thermofluid-Driven
  Transcritical Black Swans: Lagrangian Intermittency in Supercritical Fluid Turbulence}
  at the CTR Summer Program, Stanford University. \certlink{ctr}{see certificate}
  \item \textbf{2024} -- Oral presentation at ECCOMAS 2024, Lisbon (03--07/06/2024):
  \emph{Dynamics of Secondary Motions in Buoyant High-Pressure Transcritical Duct Flows}.
  \certlink{eccomas}{see certificate}
  \item \textbf{2024} -- Seminars on Computational Mechanics, EEBE--UPC: \emph{Turbulent
  square-duct flows at high-pressure transcritical conditions}.\nocert
  \item \textbf{2023} -- Open Accelerated Computing Summit (online): \emph{GPU acceleration
  of the massively-parallel flow solver RHEA with MPI+OpenACC}.\nocert
\end{itemize}

\subsection{Competitive supercomputing allocations}\label{sec:hpc}
\begin{center}\small
\begin{tabularx}{\textwidth}{@{}p{1.4cm}p{3.3cm}Xp{2.4cm}@{}}
\toprule
\textbf{Year} & \textbf{Call / system} & \textbf{Project} & \textbf{Resources} \\
\midrule
2026--27 & RES IM-2026-1-0034, MareNostrum 5 (BSC) & Extreme thermofluid events in transcritical flows: Lagrangian intermittency in supercritical fluid turbulence (PI: thesis supervisor) & 4 M core-h \\
2026 & Digital Research Alliance of Canada, RAC-2026 kfw-162-ab & GPU-accelerated simulations of microplastics transport in ocean turbulence (PI: M.\ Rahmani) & 1340 core-years \\
2025--26 & RES IM-2025-2-0017, MareNostrum 5 (BSC) & Dataset generation for turbulent flow reconstruction in external aerodynamics using sparse sensor data (PI: J.\ Calafell) & 1.7 M GPU-h \\
2025 & RES IM-2025-1-0017, MareNostrum 5 (BSC) & Same project (PI: J.\ Calafell) & 4 M GPU-h \\
2023--24 & RES IM-2023-2-0005, CTE-AMD (BSC) & Direct numerical simulations of microconfined high-pressure transcritical fluid turbulence (PI: thesis supervisor) & 490,000 core-h \\
2023 & RES IM-2023-1-0005, Finisterrae III-GPU (CESGA) & Same project & 500,000 core-h \\
\bottomrule
\end{tabularx}
\end{center}
I took part in all these allocations as a member of the research team and contributed to
writing several of the proposals.

\subsection{Research software}
\begin{itemize}
  \item \textbf{RHEA / RHEA-X} -- Open-source, GPU-resident solver (MPI + OpenACC/CUDA) for
  the high-fidelity simulation of compressible and supercritical multiphysics turbulence
  [\hyperlink{pub-J1}{J1}]. I am a co-developer of RHEA and responsible for its GPU
  acceleration (\href{https://gitlab.com/ProjectRHEA/flowsolverrhea}{gitlab.com/ProjectRHEA/flowsolverrhea}).
  \item \textbf{GPU hackathons} -- GPU porting and kernel-level optimisation of a
  thermodynamic surrogate model at the 8th MareNostrum Hackathon (BSC, 2026), and OpenACC
  acceleration of RHEA at the 6th MareNostrum Hackathon (BSC, 2024).\nocert
\end{itemize}

\subsection{Research projects}\label{sec:projects}
\begin{itemize}
  \item \textbf{SCRAMBLE} -- \emph{Turbulence-On-a-Chip: Supercritically Overcoming the
  Energy Frontier in Microfluidics}. European Research Council (ERC) Starting Grant,
  Horizon Europe, grant agreement 101040379 (UPC reference E01704). Principal
  investigator: L.\ Jofre. I am a member of the project team and have been funded by it
  as a predoctoral researcher (2023--2026) and as a postdoctoral researcher (since
  09/2026), working on the computational modelling of supercritical turbulent flows (lines
  1 and 4, Section~\ref{sec:lines}). \certlink{contract-psr}{see project addendum}
\end{itemize}

\subsection{Awards and recognition}
\begin{itemize}
  \item \textbf{2026} -- \emph{Editor's Pick} of \emph{Physics of Fluids} for the article
  [\hyperlink{pub-J2}{J2}].
  \item \textbf{2026} -- PhD thesis graded Excellent \emph{cum laude}.
\end{itemize}

% =====================================================================
\section{Internationalisation}
% \NOTE{Criterion 4 (up to 5 points).}
\subsection{International research stays}\label{sec:stays}
\begin{itemize}
  \item \textbf{28/06 -- 24/07/2026 (4 weeks, completed)} -- \textbf{Center for Turbulence
  Research, Stanford University (USA)}, 2026 Summer Program. \certlink{ctr}{see certificate}
  \begin{itemize}
    \item Project: \emph{Hunting Thermofluid-Driven Transcritical Black Swans:
    Lagrangian Intermittency in Supercritical Fluid Turbulence} (selected proposal).
    \item Topic: extreme events and causality in supercritical fluid turbulence.
    \item Team: \textbf{C.\ Monteiro}, L.\ Jofre, D.\ Mart\'in Ribes and F.\ Pizzi.
  \end{itemize}
  \item \textbf{Summer 2027, 13 weeks (scheduled)} -- \textbf{University of British
  Columbia, Vancouver (Canada)}.
  \begin{itemize}
    \item Host: \href{https://chbe.ubc.ca/mona-rahmani/}{Prof.\ Mona Rahmani},
    Department of Chemical and Biological Engineering.
    \item Topic: a causality framework to investigate the aggregation of
    microplastics in oceanic turbulent flows, building on the RAC-2026 GPU allocation on
    microplastics transport (Section~\ref{sec:hpc}).
    \item Funding: \emph{Mitacs Globalink Research Award}; the application has
    passed the internal selection of the University of British Columbia.
  \end{itemize}
  \item \textbf{2028, several months (scheduled)} -- \textbf{Universit\`a del Salento, Lecce
  (Italy)}.
  \begin{itemize}
    \item Host: group of
    \href{https://www.unisalento.it/scheda-utente/-/people/mario.direnzo}{Prof.\ Mario Di Renzo}.
    \item Topic: hypersonic flows, extreme wall-heat-flux events and causality,
    continuing the collaboration that started at the CTR 2026 Summer Program.
    \item Organisation: the Department of Fluid Mechanics has agreed to concentrate
    my teaching in one semester that year so that I can carry out the stay.
  \end{itemize}
\end{itemize}

\subsection{International training and mobility}
\begin{itemize}
  \item \textbf{2019--2020} -- Erasmus+ semester at the Politecnico di Milano (Italy).
  \certlink{erasmus}{see certificate}
  \item \textbf{2016--2026} -- Higher education in two countries: Bachelor's and Master's
  degrees at the Universidade NOVA de Lisboa (Portugal) and PhD at the UPC (Spain).
  \item \textbf{2023--2024} -- International schools and congresses: CFDParSchool at the
  Gran Sasso Science Institute (L'Aquila, Italy, 2023), machine-learning course at the von
  Karman Institute (Belgium, 2023), Compute Ontario Summer School (Canada, 2024) and
  ECCOMAS Congress (Lisbon, 2024).
\end{itemize}

\subsection{Ongoing collaborations and collaborations to be strengthened}\label{sec:collab}
\begin{itemize}
  \item \textbf{University of British Columbia (Canada)} --
  \href{https://chbe.ubc.ca/mona-rahmani/}{Prof.\ Mona Rahmani}. Causality in oceanic
  turbulent flows and aggregation of microplastics. The collaboration is already active
  through the RAC-2026 GPU allocation on microplastics transport, led by M.\ Rahmani
  (Section~\ref{sec:hpc}), and will be strengthened with the 13-week stay in summer 2027 \textbf{(scheduled)}
  and joint publications.
  \item \textbf{Universit\`a del Salento, Lecce (Italy)} --
  \href{https://www.unisalento.it/scheda-utente/-/people/mario.direnzo}{Prof.\ Mario Di Renzo}
  and \href{https://www.unisalento.it/scheda-utente/-/people/donatella.passiatore}{Dr.\ Donatella Passiatore}.
  Causal and physical analysis of extreme wall-heat-flux events in hypersonic flows. The
  collaboration started at the CTR 2026 Summer Program (Stanford) and will be strengthened
  with the research stay in 2028 \textbf{(scheduled)} and joint publications.
  \item \textbf{California Institute of Technology (USA) and Massachusetts Institute of
  Technology (USA)} --
  \href{https://www.eas.caltech.edu/people/adrian-lozano-duran}{Prof.\ Adri\'an
  Lozano-Dur\'an} (Caltech) and
  \href{https://sites.google.com/view/alvaro-martinez-sanchez/about-me}{\'Alvaro
  Mart\'inez-S\'anchez} (MIT). Causality and Lagrangian information-theoretic frameworks
  for the analysis of turbulent flows.
  \item \textbf{The Pennsylvania State University (USA)} --
  \href{https://www.me.psu.edu/department/directory-detail-g.aspx?q=xzy48}{Prof.\ Xiang
  Yang}. AI-assisted prediction of wall-normal profiles with compact neural networks
  trained on the supercritical-flow datasets developed within the ERC project SCRAMBLE
  (Section~\ref{sec:projects}).
  \item \textbf{Center for Turbulence Research, Stanford University (USA)}. Joint work in
  the CTR 2026 Summer Program project (Section~\ref{sec:stays}).
  Collaboration with Anjini Chandra (CTR), started during the 2026 Summer Program, on the
  application of non-stationary transfer-entropy frameworks to extreme wall-heat-flux events
  in supercritical fluid flows.
  \item \textbf{Stanford University (USA)} -- group of Prof.\ Matthias Ihme. Contribution
  to the BLASTNet multiphysics simulation dataset for machine learning [\hyperlink{pub-D1}{D1}].
  \item \MAYBE{\textbf{Smart Ocean Peniche (Portugal)} -- Sea Science and Technology Park
  of Peniche, an ocean-economy innovation ecosystem promoted by the Municipality of
  Peniche, Docapesca, the Polytechnic Institute of Leiria and Biocant (prospective).
  Support to the R\&D projects of its companies in ocean engineering, extending to applied
  problems the expertise in oceanic turbulent flows developed with UBC (Canada).}
  (\href{https://smartoceanpeniche.com}{smartoceanpeniche.com})
  \item \MAYBE{\textbf{Instituto Superior T{\'e}cnico, Lisbon (Portugal)} -- Prof. Carlos da Silva
  (prospective). Turbulent/turbulent (TT) and turbulent/non-turbulent (TNT) interfaces in
  flows with variable fluid properties, extending the supercritical-fluid DNS database.}
\end{itemize}

% =====================================================================
\section{Professional experience and technology transfer}\label{sec:prof}

\subsection{Thermal simulation of automotive electronics and electric drives}
Continental Engineering Services, Porto (Portugal), 2021--2022. Role: Thermal Simulation
Engineer. \certlink{continental}{see certificate}
\begin{itemize}
  \item \textbf{Thermal reliability of power electronics.}
  \begin{itemize}
    \item Thermal-management and thermal-reliability simulations of electronic components
    and inverters.
    \item Verification of operating temperatures and design margins of automotive control
    units.
    \item Tools: Ansys Icepak.
  \end{itemize}
  \item \textbf{Cooling of electric motors.}
  \begin{itemize}
    \item Conjugate heat-transfer CFD simulations of electric-motor cooling systems, in
    support of their thermal design.
    \item Tools: STAR-CCM+ and Ansys.
  \end{itemize}
  \item \textbf{Tolerance stack-up analysis.}
  \begin{itemize}
    \item Dimensional tolerance (stack-up) calculations of assemblies, carried out
    together with the mechanical design team.
  \end{itemize}
\end{itemize}

\subsection{Vehicle aerodynamics}
\begin{itemize}
  \item \textbf{Aerodynamic drag reduction of trucks} (Continental Engineering Services,
  internal project, 2021--2022). Role: team member.
  \begin{itemize}
    \item CFD study of the external aerodynamics of heavy vehicles.
    \item Evaluation of drag-reduction solutions.
  \end{itemize}
  \item \textbf{Aerodynamic package of the FST~10e Formula Student car} (FST Lisboa,
  Universidade de Lisboa, 2020--2021; non-salaried student team). Role: member of the
  Aerodynamics department. (\href{http://hdl.handle.net/10362/149236}{see Master's thesis})\nocert
  \begin{itemize}
    \item Simulation: CFD of the external aerodynamics of the car.
    \item Geometry: preparation and cleaning of the geometry for meshing.
    \item Design: multi-element rear wing and its drag-reduction system (DRS).
    \item Manufacturing: carbon-fibre composite parts.
    \item Outcome: Master's thesis (graded 19/20) and the SAE paper [\hyperlink{pub-C2}{C2}].
  \end{itemize}
\end{itemize}

\subsection{Design engineering and transfer of research software and data}
\begin{itemize}
  \item \textbf{Production-oriented mechanical design in the metalworking industry}
  (Cubot\'onic -- Ind\'ustria Metalomec\^anica, Malveira, Portugal, 2019). Role: Design
  Engineering Intern.\nocert
  \begin{itemize}
    \item Products: custom cabinets, machines and metal structures, designed in the
    company's engineering office from the client's requirements to the manufacturing
    drawings.
    \item Design for manufacturing: every design had to suit the in-house production
    chain (laser cutting and punching, bending and clinching, machining, welding and
    metalwork, powder coating and final assembly), which meant choosing sheet thicknesses,
    bend radii, joints and tolerances that could be manufactured efficiently and assembled
    without rework.
    \item Tools: Siemens Solid Edge.
  \end{itemize}
  \item \textbf{Open-source simulation software and open data} (UPC, 2023--2026). Role:
  co-developer of RHEA, responsible for its GPU acceleration, and co-author of the databases.
  \begin{itemize}
    \item Release of the GPU-resident solver RHEA-X [\hyperlink{pub-J1}{J1}] as open-source software.
    \item Open high-fidelity databases of microconfined transcritical flows [\hyperlink{pub-J3}{J3}, \hyperlink{pub-D2}{D2}].
    \item Contribution to the BLASTNet machine-learning dataset [\hyperlink{pub-D1}{D1}].
    \item Impact: results reusable by academia and industry.
  \end{itemize}
\end{itemize}

\subsection{Technical skills}
\begin{itemize}
  \item \textbf{Programming and HPC} -- C++, Python, Julia, MATLAB; CUDA, OpenACC and MPI;
  Linux and Slurm HPC environments.
  \item \textbf{Simulation} -- STAR-CCM+, Ansys (Fluent, CFX, Icepak).
  \item \textbf{Computer-aided design} -- SolidWorks, Siemens NX, Siemens Solid Edge and
  Creo.
\end{itemize}

% =====================================================================
\section{Adequacy to the position profile}\label{sec:adequacy}
The position LE-403/729 combines three requirements: teaching in the thermofluid subjects
of the EEBE and the Master's Degree in Mechanical Technologies; research in the three lines
of the Department of Fluid Mechanics; and the supervision of doctoral theses. My
trajectory addresses all three. It combines a mechanical-engineering education with a
strong thermofluid component, industrial experience in thermal and aerodynamic
simulation, and a PhD and postdoctoral activity in the department on the high-fidelity
simulation of supercritical turbulence. Each requirement is examined below in the order of
the call, from three perspectives: the background that supports it, the current activity
that demonstrates it, and the projection of my career, as detailed in the Work Plan, that
will consolidate it.

\subsection{Teaching profile}
I have taught subjects of the profile at the EEBE and the EPSEVG since 2024--25, with 180
contact hours up to 2026--27 S1, a full load of 242 hours in 2026--27, and a certified
student rating of 4.42 out of 5 in 2025--26 (Section~\ref{sec:surveys}).
\begin{itemize}
  \item \textbf{Thermodynamics and Heat Transfer; Thermodynamics.} My education in heat
  transfer, the thermal simulation of electronics and electric motors at Continental, and
  a PhD centred on real-fluid thermodynamics and heat transfer provide the basis for these
  subjects. In 2026--27 I teach the full theory and problem course of group M1 and support the
  coordination, examinations and marking of the subject, with a further group in the
  spring semester \textbf{(scheduled)}. I intend to integrate real-fluid property tools and
  computational examples into the course and to assume coordination tasks progressively.
  \item \textbf{Fluid Mechanics.} The subject builds directly on my Master's courses in
  Computational Fluid Mechanics (18/20) and Aerodynamics (20/20) and on my doctoral work
  on wall-bounded turbulence. I have taught its laboratory every year since 2024--25 (eight
  groups, 120 hours), returning individual feedback on every report. My aim is to extend
  this activity to theory and problem sessions, and to connect the laboratory with
  simulation through computational notebooks.
  \item \textbf{Fluid Engineering.} The aerodynamic design of a Formula Student car,
  documented in my Master's thesis and in [\hyperlink{pub-C2}{C2}], and the drag reduction
  of trucks at Continental constitute my applied background in external flows, now
  complemented by a DNS database of a NACA~0012 airfoil [\hyperlink{pub-R5}{R5}]. I will
  teach the theory of groups M65 and M66 at the EPSEVG in spring 2027 (92 hours)
  \textbf{(scheduled)}, following the course development of the Work Plan, with applied
  cases in external aerodynamics and turbomachinery.
  \item \textbf{Thermal Engineering; Power Plants and Renewable Energies;
  Thermofluid-Dynamic Generation.} Industrial thermal design and research on supercritical
  fluids, the working media of sCO$_2$ power cycles, support these subjects. I currently
  teach the laboratory of Fundamentals of Thermal Engineering at the EPSEVG and develop
  real-fluid surrogate models [\hyperlink{pub-R4}{R4}]. In the medium term, I intend to
  introduce supercritical cycles and real-fluid effects into the thermal-power subjects.
  \item \textbf{Computational Fluid Mechanics (Master's Degree in Mechanical
  Technologies).} Industrial CFD with Ansys and STAR-CCM+ and specialised schools on CFD,
  HPC and machine learning (Sections~\ref{sec:continuing} and~\ref{sec:training}) form my
  background. I am a co-developer of the solver RHEA, responsible for its GPU acceleration
  [\hyperlink{pub-J1}{J1}], and have taken part in six competitive supercomputing
  allocations. This experience will translate into GPU laboratory sessions, data-driven
  and causal flow analysis, and course projects that lead to master's theses.
\end{itemize}

\subsection{Research profile}
My research lies at the intersection of the three lines of the call: the high-fidelity
simulation of supercritical multiphysics turbulence, the analysis of its extreme events
with causal and data-driven methods, and the GPU-accelerated computation that makes such
simulations possible. It has produced seven articles in Q1 journals (2024--2026), four of
them as first author, and four further submitted articles
(Section~\ref{sec:publications}).
\begin{itemize}
  \item \textbf{Multiphysics turbulent flows; compressible and supercritical fluids.} My
  PhD thesis (Excellent \emph{cum laude}) established resolution standards, minimal flow
  units and flow regimes of transcritical wall turbulence
  [\hyperlink{pub-J3}{J3}--\hyperlink{pub-J7}{J7}]. As a postdoctoral researcher in the
  ERC project SCRAMBLE, I currently study buoyancy and secondary flows in square ducts
  [\hyperlink{pub-J2}{J2}, \hyperlink{pub-R1}{R1}], work selected as an \emph{Editor's
  Pick} and to be presented at EFDC3 in 2027 \textbf{(scheduled)}. The Work Plan extends
  this line towards high-fidelity databases and the analysis of the coupling between
  turbulence, heat transfer and real-fluid thermodynamics (WP1).
  \item \textbf{Data-driven prediction of engineering flows; extreme events; causality.}
  Open databases for machine learning [\hyperlink{pub-J3}{J3}, \hyperlink{pub-D1}{D1},
  \hyperlink{pub-D2}{D2}] and the data-informed analysis of experiments
  [\hyperlink{pub-J7}{J7}] laid the foundation of this line. Since the CTR 2026 Summer
  Program, I have applied it to extreme wall-heat-flux events
  [\hyperlink{pub-R2}{R2}, \hyperlink{pub-R3}{R3}], in collaboration with CTR, Caltech, MIT
  and Penn State, and will present it at THMT-27 \textbf{(scheduled)}. Its projection is
  the identification of causal mechanisms and the causality-informed prediction of extreme
  events (WP2--WP3), supported by the stays at UBC in 2027 and at the Universit\`a del
  Salento in 2028 \textbf{(scheduled)}.
  \item \textbf{Modelling and computation; distributed (CPU) and accelerated (GPU)
  programming.} Industrial CFD and GPU porting with OpenACC since 2023 (Open Accelerated
  Computing Summit, BSC Hackathons) form the background of this line. My current work
  comprises the GPU acceleration of RHEA [\hyperlink{pub-J1}{J1}], low-rank thermodynamic
  surrogates [\hyperlink{pub-R4}{R4}] and six supercomputing allocations totalling more
  than 10 million core and GPU hours. The Work Plan projects it towards
  performance-portable open-source solvers for pre-exascale systems and real-fluid
  turbulence models (WP4).
\end{itemize}

\subsection{Supervision and integration in the department}
\begin{itemize}
  \item \textbf{Supervision of theses.} I co-supervised a Bachelor's thesis (2023),
  currently tutor a curricular internship on extreme events in turbulent flows, and will
  serve on a TFG assessment committee in 2027 \textbf{(scheduled)}. Building on this
  experience, the Work Plan proposes the co-supervision of two doctoral theses during the
  contract, on topics drawn from the three research lines, together with TFG and TFM
  projects that feed these lines.
  \item \textbf{Integration in the department.} I have been a member of the research group
  GReCEF since 2023, I am funded by the ERC project SCRAMBLE, and I already share the
  teaching of Thermodynamics and Heat Transfer with the lecturers of the department. I
  teach in Spanish, use English as the working language of my research, and am committed
  to reaching the C1 level of Catalan (Section~\ref{sec:languages}).
\end{itemize}

In summary, my background, my current activity and the projection set out in the Work
Plan converge on the teaching and research needs of the position. I therefore consider
myself prepared to take on the responsibilities of the Lector position in the Department
of Fluid Mechanics of the Universitat Polit\`ecnica de Catalunya, linked to the Barcelona
East School of Engineering (EEBE).

% =====================================================================
\newpage
\section{Supporting documents}
List of attached documents, in the order in which they are cited:
\begin{enumerate}[itemsep=0pt]
\item Academic qualifications: PhD degree fee receipt (UPC); Master's diploma
(\emph{Carta de Curso}); Master's and Bachelor's certificates, Diploma Supplements and Spanish grade equivalences (NOVA Lisbon);
Erasmus+ certificate (Politecnico di Milano).
\item Contracts (personal data redacted): substitute professor, postdoctoral researcher
(with ERC addendum) and predoctoral researcher; CTR-Stanford Summer Program letter;
Continental work certificate.
\item Teaching: UPC certificate of teaching activity (hours, TFG supervision and student surveys) and 2026--27 teaching assignment; TTC statement by
the course coordinator; internship training project; detailed student surveys; Catalan certificates (SLT-UPC).
\item Training: ICE course; CFD, HPC and machine-learning schools and courses.
\item Research: ECCOMAS attendance
certificate.
\end{enumerate}

\certpdf{phd}{01_education/PhD_UPC_2026_degree_fee_receipt.pdf}
\certpdf{msc-dip}{01_education/MSc_NOVA_2021_diploma.pdf}
\certpdf{msc}{01_education/MSc_NOVA_2021_certificate.pdf}
\certpdf{msc-sup}{01_education/MSc_NOVA_2021_diploma_supplement.pdf}
\certpdf{msc-eq}{01_education/MSc_NOVA_2021_grade_equivalence.pdf}
\certpdf{erasmus}{01_education/Erasmus_Polimi_2019-20_certificate.pdf}
\certpdf{bsc}{01_education/BSc_NOVA_2019_certificate.pdf}
\certpdf{bsc-sup}{01_education/BSc_NOVA_2019_diploma_supplement.pdf}
\certpdf{bsc-eq}{01_education/BSc_NOVA_2019_grade_equivalence.pdf}
\certpdf{contract-sub}{02_positions/Contract_substitute_lecturer_2026-27_REDACTED.pdf}
\certpdf{contract-psr}{02_positions/Contract_postdoc_SCRAMBLE_2026_REDACTED.pdf}
\certpdf{contract-pre}{02_positions/Contract_predoctoral_2023-2026_REDACTED.pdf}
\certpdf{ctr}{06_international/CTR_Stanford_Summer_Program_2026_letter.pdf}
\certpdf{continental}{02_positions/Continental_work_certificate.pdf}
\certpdf{teaching}{03_teaching/UPC_teaching_certificate_2026_REDACTED.pdf}
\certpdf{assignment}{03_teaching/Teaching_assignment_2026-27.pdf}
\certpdf{ttc-coord}{03_teaching/TTC_coordination_statement_Jofre_2026_signed.pdf}
\certpdf{internship}{03_teaching/Internship_605632_training_project_REDACTED.pdf}
\certpdf{surveys}{03_teaching/Student_surveys_MF_2025-26_S2.pdf}
\certpdf{ice1}{04_training/ICE_M2-4_Atenea_2026_certificate.pdf}
\certpdf{monsha}{04_training/Monsha_AI_Educator_L1_certificate.pdf}
\certpdf{cfdpar}{04_training/CFDParSchool_2023_admission_letter.pdf}
\certpdf{vki}{04_training/VKI_ML_Fluid_Dynamics_2023_certificate.pdf}
\certpdf{python}{04_training/Flowthermolab_Python_for_Engineers_2022_certificate.png}
\certpdf{ml-ftl}{04_training/Flowthermolab_ML_Mechanical_Engineers_2023_certificate.pdf}
\certpdf{ansys}{04_training/Ansys_Turbulence_Beyond_Basics_2022_certificate.pdf}
\certpdf{ansys-srs}{04_training/Ansys_Scale_Resolving_Simulations_2022_certificate.pdf}
\certpdf{esss}{04_training/ESSS_Ansys_Workshop_2020_certificate.pdf}
\certpdf{simscape}{04_training/MathWorks_Simscape_Onramp_2021_certificate.pdf}
\certpdf{simulink}{04_training/MathWorks_Simulink_Onramp_2021_certificate.pdf}
\certpdf{udemy3}{04_training/Udemy_CFD_Fundamentals_3_2021_certificate.pdf}
\certpdf{udemy}{04_training/Udemy_CFD_Fundamentals_2_2020_certificate.pdf}
\certpdf{udemy1}{04_training/Udemy_CFD_Fundamentals_1_2020_certificate.pdf}
\certpdf{eccomas}{06_international/ECCOMAS_2024_attendance_certificate.pdf}
\certpdf{catalan1}{07_other/Catalan_A2.1_SLT-UPC_2025_certificate.pdf}
\certpdf{catalan2}{07_other/Catalan_A2.2_SLT-UPC_2025_certificate.pdf}

\end{document}
